\begin{textsample}{POS Dim 1 – 2016\_c4 – Score -4.00 – 2016\_c4/enade\_2016\_tec\_rad\_371c7e6d01.txt}  \label{ex:f1_pos_288}
CAPÍTULO 4 – PERCEPÇÃO DA PROVA As análises feitas neste capítulo tratam das percepções dos concluintes da Área de Tecnologia em Radiologia sobre a prova aplicada no Enade/2016.
Estas percepções foram mensuradas por meio de nove questões que avaliaram desde o grau de dificuldade da prova até o tempo gasto para concluí -la.
As percepções sobre a prova foram relacionadas com o desempenho dos estudantes e com a Grande Região de funcionamento do curso.
O questionário de percepção da prova encontra -se ao final do Anexo VII, que traz a reprodução da prova.
O desempenho dos estudantes foi classificado em quatro quartos.
Para tanto, esse desempenho foi ordenado de forma ascendente.
O percentil 25, P25, também conhecido como primeiro quartil, é a nota de desempenho que deixa um quarto ( 25 \% ) dos valores observados abaixo e três quartos acima.
A Figura 1 apresenta uma ilustração deste conceito.
O quarto inferior de desempenho é composto pelas notas abaixo do primeiro quartil.
Já o percentil 75, P75, também conhecido como terceiro quartil, é o valor para o qual há três quartos ( 75 \% ) dos dados abaixo e um quarto acima dele.
O quarto superior de desempenho é composto pelas notas iguais ou acima do terceiro quartil.
O percentil 50, P50, também conhecido como mediana, é o valor que divide as notas em dois conjuntos de igual tamanho.
O segundo quarto inclui valores entre o primeiro quartil ( P25 ) e a mediana.
O terceiro quarto contém os valores entre a mediana ( P50 ) e o terceiro quartil ( P75 ).
Vale ressaltar que percentis, quartis e medianas são pontos que não obrigatoriamente pertencem ao conjunto original de dados, ao passo que os quartos são subconjuntos dos dados originais.
A seguir, serão apresentados gráficos com resultados selecionados, relativos às nove questões avaliadas por grupos de estudantes.
Os gráficos apresentam nas barras o percentual de alunos que assinalaram uma das opções ou a soma das porcentagens daqueles que assinalaram duas ( ou três ) delas.
Por exemplo, para as questões 1 e 2, os gráficos apresentam a porcentagem total de participantes que assinalaram as opções ( D ) Difícil e ( E ) Muito difícil.
Em cada barra foram assinalados também os extremos do intervalo de confiança de 95 \% como linhas verticais unidas por uma linha horizontal na forma da letra H maiúscula.
As Tabelas no Anexo II apresentam os valores absolutos e a distribuição percentual18 das alternativas válidas das nove questões, segundo o mesmo recorte de desempenho dos alunos e Grande Região de funcionamento do curso. 4.1 - GRAU DE DIFICULDADE DA PROVA 4.1.1 - Componente de Formação Geral Ao avaliarem “ Qual o grau de dificuldade desta prova na parte de Formação Geral? ” ( Questão 1 ), 31,3 \% do grupo de inscritos e presentes optaram pelas alternativas Difícil ou Muito difícil.
Entretanto, para mais da metade dos estudantes ( 62,1 \% ), o Componente de Formação Geral da prova foi considerado com grau de dificuldade Médio ( Gráfico 4.1, Gráfico 4.2 e, no Anexo II, a Tabela II.1 ).
O percentual de estudantes que consideraram a prova como Difícil ou Muito difícil foi maior na região Sudeste, onde a proporção foi de 34,3 \%, enquanto a de menor incidência foi a Nordeste, com 25,5 \%.
No Gráfico 4.1, é possível observar que a diferença entre a região Nordeste e a região Sudeste é estatisticamente significativa.
Nas Grandes Regiões, a proporção de presentes à prova que consideraram o Componente de Formação Geral como sendo de grau de dificuldade Médio esteve entre 59,4 \% na região Sul e 66,4 \% na região Nordeste. 18 Cumpre lembrar uma das convenções para tabelas numéricas ( pág. iii ) sobre a possibilidade de a soma das partes não resultar em 100 \% por questões de arredondamento.
O percentual de alunos que consideraram a prova Difícil ou Muito difícil foi decrescente em função dos quartos de desempenho : 35,2 \% no primeiro quarto e 26,4 \% no quarto, grupo de melhor desempenho na prova.
Nos quartos de desempenho intermediários, a proporção de alunos que consideraram a prova Difícil ou Muito difícil foi de 33,0 \% no segundo quarto e 30,6 \% no terceiro quarto.
A diferença entre o primeiro e último quartos é estatisticamente significativa.
Para todos os quartos de desempenho, a alternativa modal para esta pergunta foi Médio, com 59,9 \% e 62,7 \% dos respondentes nos quartos extremos, primeiro e quarto, respectivamente. 4.1.2 - Componente de Conhecimento Específico Ao responderem à Questão 2 – “ Qual o grau de dificuldade desta prova na parte de Componente Específico? ” – 42,1 \% do grupo de estudantes classificaram -na como Difícil ou Muito difícil.
Além disso, o Componente de Conhecimento Específico da prova foi considerado com grau de dificuldade Médio por 54,4 \% dos alunos ( Gráfico 4.3, Gráfico 4.4, e, no Anexo II, a Tabela II.2 ).
A análise das respostas dos estudantes quanto ao grau de dificuldade do Componente de Conhecimento Específico da prova, agregado por Grande Região, mostra que as diferenças entre a maior e as duas menores proporções de alunos que a avaliaram como Difícil ou Muito difícil são estatisticamente significativas : a maior na região Sul ( 47,4 \% ) e as menores proporções na região Norte ( 35,1 \% ), seguida da região Nordeste ( 37,0 ).
O percentual de alunos que classificaram o grau de dificuldade como Médio, no Componente de Conhecimento Específico, variou de 49,5 \% a 61,7 \%, para as regiões Sul e Norte, respectivamente.
Considerando -se a avaliação da dificuldade das questões do Componente de Conhecimento Específico da prova, de acordo com o desempenho dos estudantes, observa -se que não há diferença estatisticamente significativa dos resultados entre todos os quartos de desempenho.
A proporção dos que classificaram a parte específica como Difícil ou Muito difícil variou de 39,3 \% ( último quarto ) a 42,4 \% ( primeiro quarto ).
A alternativa modal para a Questão 2 foi o grau Médio, com 53,4 \% do quarto inferior e 56,4 \% do superior optando por esta resposta. 4.2 - EXTENSÃO DA PROVA EM RELAÇÃO AO TEMPO TOTAL Indagados quanto à extensão da prova, em relação ao tempo total oferecido para a sua resolução ( Questão 3 ), os estudantes apontaram, com maior incidência, a alternativa que considerava a extensão adequada, para todas as agregações consideradas ( Gráfico 4.5, Gráfico 4.6, e, no Anexo II, a Tabela II.3 ).
O percentual de alunos que responderam ser a extensão da prova adequada foi de 70,5 \%.
Já 24,0 \% dos inscritos presentes consideraram que a prova foi longa ou muito longa, e 5,4 \% a avaliaram como curta ou muito curta.
Entre as Grandes Regiões, a proporção daqueles que avaliaram a prova como longa ou muito longa em relação ao tempo total destinado à sua resolução variou de 12,6 \% na região Norte até 27,1 \% na região Sudeste.
A diferença entre a região Norte e as demais é estatisticamente significativa.
Considerando -se o desempenho dos alunos, nota -se ainda que 68,0 \% consideraram a extensão da prova adequada no quarto de desempenho inferior e 72,5 \% no de melhor desempenho ( quarto superior ).
Nos quartos intermediários, esta proporção foi 71,3 \% no segundo quarto e 70,3 \% no terceiro.
No Gráfico 4.6, pode -se constatar que há uma tendência decrescente da proporção de estudantes que consideraram a prova longa ou muito longa em função dos quartos de desempenho.
Observa -se que não há diferença estatisticamente significativa entre as proporções de escolha desta resposta por quartos de desempenho. 4.3 - COMPREENSÃO DOS ENUNCIADOS DAS QUESTÕES 4.3.1 - Componente de Formação Geral Com relação aos enunciados das questões do Componente de Formação Geral ( Questão 4 ), as opiniões foram positivas, já que 77,1 \% dos alunos avaliados consideraram todos ou a maioria dos enunciados das questões “ claros e objetivos ” ( Gráfico 4.7, Gráfico 4.8, e, no Anexo II, a Tabela II.4 ).
Na análise regional, a percentagem de estudantes que avaliaram que todos ou a maioria dos enunciados das questões do Componente de Formação Geral estavam “ claros e objetivos ” variou de 73,0 \% na região Sudeste a 83,1 \% na região Centro-Oeste, sendo esta diferença estatisticamente significativa.
A análise das percepções dos estudantes sobre a clareza e objetividade dos enunciados permite afirmar que todos, ou a maioria dos enunciados de questões relativas ao Componente de Formação Geral, foram considerados “ claros e objetivos ” para a maior parte dos respondentes ( igual ou maior do que 73 \% em todas as regiões e igual ou maior do que 71,2 \% para todos os quartos de desempenho ).
Segundo o desempenho, observa -se que a proporção dos que emitiram esta opinião cresce conforme o desempenho aumenta, com diferenças estatisticamente significativas entre o primeiro e os dois últimos quartos de desempenho.
No quarto superior, a clareza e objetividade de todos ou da maioria dos enunciados das questões foi percebida por 80,4 \%. 4.3.2 - Componente de Conhecimento Específico Com relação aos enunciados das questões do Componente de Conhecimento Específico da prova, para 80,3 \% dos estudantes avaliados da Área de Tecnologia em Radiologia, a clareza e a objetividade ( Questão 5 ) estavam presentes em todas ou na maioria das questões ( Gráfico 4.9, Gráfico 4.10, e no Anexo II, a Tabela II.5 ).
A maioria dos estudantes de todas as Grandes Regiões brasileiras considerou “ claros e objetivos ” todos ou a maioria dos enunciados das questões do Componente de Conhecimento Específico da prova, percentual sempre maior ou igual a 78,7 \%.
As diferenças entre as Grandes Regiões não são estatisticamente significativas.
A proporção de estudantes que consideraram os enunciados das questões “ claros e objetivos ” apresenta uma tendência crescente em relação ao aumento de desempenho : mais elevada no quarto superior ( 85,0 \% ) se comparada ao quarto inferior de desempenho ( 74,4 \% ).
As diferenças entre o primeiro e os dois últimos quartos de desempenho, e do segundo para o quarto de melhor desempenho são estatisticamente significativas. 4.4 - SUFICIÊNCIA DAS INFORMAÇÕES/INSTRUÇÕES FORNECIDAS Ao avaliarem as informações/instruções fornecidas para a resolução das questões ( Questão 6 ), 82,4 \% dos respondentes da Área de Tecnologia em Radiologia de todo o Brasil afirmaram que estas eram até excessivas ou “ suficientes ” em todas ou na maioria das questões ( Gráfico 4.11, Gráfico 4.12, e, no Anexo II, a Tabela II.6 ).
Quanto à distribuição de respondentes pelas Grandes Regiões observa -se que a proporção de estudantes que consideraram as informações/instruções fornecidas até excessivas ou “ suficientes ” em todas ou na maioria das questões foi sempre igual ou superior a 81,3 \%, chegando a 83,7 \% na região Centro-Oeste.
A diferença entre o maior percentual, obtido na região Centro-Oeste ( 83,7 \% ), e o menor, da região Norte ( 81,3 \% ) não é estatisticamente significativa.
Levando -se em conta o desempenho dos participantes, notam -se diferenças estatisticamente significativas entre as opiniões de estudantes de alguns quartos de desempenho, como mostra o Gráfico 4.12.
O percentual de participantes que avaliaram “ as informações/instruções ” como até excessivas ou “ suficientes ” em todas ou na maioria das questões foi mais elevado no quarto superior ( 88,5 \% ), percentual superior à média nacional ( 82,4 \% ).
Já no quarto inferior, a suficiência das informações/instruções declarada como até excessiva, em todas ou na maioria das questões foi percebida por 74,3 \% dos respondentes.
Estas proporções são crescentes com os quartos de desempenho. 4.5 - DIFICULDADE ENCONTRADA AO RESPONDER À PROVA Perguntados sobre as dificuldades com as quais depararam ao responder à prova ( Questão 7 ), 13,1 \% dos estudantes apontaram o Desconhecimento do conteúdo.
Para 54,9 \%, a Forma diferente de abordagem do conteúdo foi indicada como dificuldade.
Já a Falta de motivação para fazer a prova foi a dificuldade apontada por 8,7 \% dos respondentes.
Considerando -se todo o Brasil, 20,5 \% dos respondentes afirmaram que não tiveram qualquer tipo de dificuldade para responder à prova ( Tabela II.7 no Anexo II ).
Os Gráficos 4.13 e 4.14 apresentam os percentuais de estudantes que apontaram o Desconhecimento do conteúdo como dificuldade percebida ao responder à prova.
Na análise por Grandes Regiões, o percentual de inscritos e presentes que apontaram o Desconhecimento do conteúdo como dificuldade ao responder à prova não superou 17,6 \%.
Os percentuais variaram de 7,7 \% na região Centro-Oeste a 17,6 \% na Sul, sendo esta diferença significativa estatisticamente.
A Forma diferente de abordagem do conteúdo foi a escolha modal dos estudantes, com percentuais que variaram de 47,3 \% ( região Nordeste ) a 60,5 \% ( Norte ).
O percentual de alunos que citaram a Falta de motivação como dificuldade variou de 6,2 \% ( região Norte ) a 11,8 \% ( região Sul ).
Os que declararam não ter qualquer tipo de dificuldade para responder à prova variaram de 18,2 \% na região Sudeste a 26,2 \% na Nordeste.
Com relação aos quartos de desempenho, o Desconhecimento do conteúdo foi a opção escolhida por 13,1 \% dos estudantes do quarto superior e 12,1 \% do quarto inferior.
A alternativa modal para os alunos, quando agregados pelos quartos de desempenho, foi que a dificuldade encontrada foi causada pela Forma diferente de abordagem do conteúdo : 53,6 \% no quarto inferior e 52,7 \% do quarto superior assim o responderam. 4.6 - CONTEÚDOS DAS QUESTÕES OBJETIVAS DA PROVA Ao analisar os conteúdos das questões objetivas da prova ( Questão 8 ), um percentual muito pequeno dos estudantes avaliados, apenas 4,1 \%, afirmou que não estudou ainda a maioria desses conteúdos ( Gráficos 4.15, Gráfico 4.16, e a Tabelas II.8 no Anexo II ).
A maioria ( 64,6 \% ) afirmou ter estudado e aprendido muitos ou todos os conteúdos avaliados.
Na análise por Grande Região, a proporção de respondentes que escolheram a opção não estudou ainda a maioria desses conteúdos foi pequena.
Observa -se que nas regiões Norte ( 7,0 \% ) e Nordeste ( 4,9 \% ), apesar de pequenas, as proporções foram maiores do que a média nacional ( 4,1 \% ).
Não foi observada diferença estatisticamente significativa entre as regiões.
Em todas as regiões, a maioria dos presentes afirmou ter estudado e aprendido muitos ou todos os conteúdos, com proporções variando entre 57,9 \% na região Norte e 68,1 \% na Centro-Oeste.
Considerando -se separadamente as opiniões de estudantes dos quatro quartos de desempenho, observa -se que, no quarto inferior, 6,7 \% ofereceram como resposta que não estudou ainda a maioria desses conteúdos, sendo 1,4 \% os do quarto superior com a mesma resposta.
As diferenças entre os alunos que optaram por este motivo de dificuldade no quarto de melhor desempenho e os dois primeiros quartos são estatisticamente significativas.
Tendo em conta o quarto superior, 78,7 \% dos alunos afirmaram ter estudado e aprendido muitos ou todos os conteúdos.
No outro extremo, no primeiro quarto, 50,7 \% optaram pelas mesmas categorias.
As proporções são crescentes com o desempenho. 4.7 - TEMPO GASTO PARA CONCLUIR A PROVA Ao responderem sobre o tempo de conclusão da prova ( Questão 9 ), a grande maioria dos estudantes ( 80,6 \% ) afirmou ter gasto Entre duas e quatro horas ( Gráfico 4.17, Gráfico 4.18 e, no Anexo II, a Tabela II.9 ).
Considerando -se as cinco Grandes Regiões brasileiras, o percentual dos que utilizaram Entre duas e quatro horas para finalizar a prova nas regiões Norte ( 74,9 \% ), Nordeste ( 79,2 \% ) e Sul ( 72,2 \% ) foi inferior ao percentual nacional.
Nas demais Grandes Regiões, o percentual de alunos que dispensaram Entre duas e quatro horas para concluir a prova foi de 84,6 \%, como mostra o Gráfico 4.17.
Uma vez sendo analisadas as alternativas escolhidas pelos estudantes que se situam nos diferentes quartos de desempenho, observa -se uma tendência crescente com o desempenho : uma maior proporção de participantes no quarto superior declarou ter gasto Entre duas e quatro horas “ para concluir a prova ” quando comparadas com os dos quartos inferiores.
As diferenças entre o primeiro ( 73,8 \% ) e o terceiro ( 82,3 \% ) e o quarto de maior desempenho ( 86,5 \% ) são estatisticamente significativas.

% matched lemmas: 
\end{textsample}
